\section{Formal Execution Model}
\label{sec:model}

To resolve the five failure modes identified in Section~\ref{sec:problem}, we formalize an execution model built around Intent, Custody, Trajectory, and Proof. This model abstracts agentic software engineering as a state transition system operating over a repository space.

Let $R$ represent the state of the repository, including the files, directory structure, git history, and runtime environment. An agent run is a sequence of state transitions $R_0 \xrightarrow{a_1} R_1 \xrightarrow{a_2} \dots \xrightarrow{a_n} R_n$, where each action $a_i$ is executed by the agent via tools.

\subsection{Intent: The Goal Objective}
We define Intent ($I$) as a durable task contract containing goals, constraints, unresolved questions, stop conditions, and proof expectations:
$$I = \langle G, C, V \rangle$$
where $G = \{g_1, g_2, \dots, g_k\}$ represents the goals, $C$ represents explicit constraints, and $V = \{v_1, v_2, \dots, v_m\}$ represents proof hooks or validation expectations. In Decapod, these concerns are distributed across governed plan artifacts, todo records, project specs, and proof configuration rather than a single universal \texttt{intent.json}. The contract persists outside the model context, allowing later steps to re-resolve intent and detect unresolved questions before execution.

\subsection{Custody: Environment Bounding}
Custody ($K$) defines the execution and security sandbox for the agent. Formally, it is a tuple:
$$K = \langle W, B, A, P \rangle$$
where:
\begin{itemize}
    \item $W \subset R$ is the isolated workspace slice (a task-scoped git worktree and, when configured, a container).
    \item $B$ is the task-specific branch locked for the agent.
    \item $A$ is the set of governance operations and proof capabilities available to the run.
    \item $P$ represents the repository and runtime boundary enforced by the workspace configuration.
\end{itemize}
By enforcing custody boundaries, the kernel keeps the agent's repository mutations in $W$ and prevents direct work on protected branches. This reduces cross-task workspace collisions; it is not a complete operating-system sandbox or credential-isolation proof.

\subsection{Trajectory: The Governance Record}
The Trajectory ($T$) is an event-backed sequence of governance records:
$$T = [t_1, t_2, \dots, t_n]$$
where a record may include timestamp, operation, actor, session, correlation identifier, status, and structured details. Decapod journals todo, broker, proof, and related governance events and can render them through its flight recorder. These records improve reconstruction and review, but they are not asserted here to be a complete capture of every model token, shell byte, file diff, or tamper-proof execution event.

\subsection{Proof: Completion Evidence}
A Proof ($P_f$) is the set of machine-checkable evidence attached to a governed task:
$$P_f = \langle I, R_n, \{\operatorname{Pass}(v_j)\}, E \rangle$$
where $E$ contains relevant proof events, workunit results, validation output, and provenance references. In Decapod, validation and proof gates can block promotion or completion paths when required evidence is absent. A proof record establishes what the configured gates observed; it does not establish semantic correctness beyond those gates, and this paper makes no claim that the current implementation signs a universal certificate.
